\subsubsection{3.2 研究内容}

本项目围绕海杂波场景电磁频谱数据智能处理机理与方法展开研究，针对强海杂波干扰下微弱电磁频谱信号提取难、有限算力端侧海量频谱数据轻量可信管理难、动态对抗环境下边缘自主认知与多平台协同可信决策难等问题，旨在构建从海杂波场景频谱信号传播与提取、端侧频谱数据轻量可信管理、边缘频谱智能处理到原型系统验证的完整理论与方法体系。首先，研究\textbf{海杂波干扰下微弱电磁频谱信号的传播规律与提取方法}，揭示复杂海况下信号时空耦合传播机理并构建干扰对消增强机制，解决强杂波背景下微弱目标“看不清、提不准”的难题，为后续处理提供高质量频谱特征输入；其次，研究\textbf{有限算力端侧频谱数据的轻量存储与可信管理方法}，通过事件驱动价值评估、语义压缩与链上链下协同存证，突破端侧资源受限下海量数据“存不下、查不快、信不过”的瓶颈，为频谱数据全生命周期提供可靠数据基座；接着，研究\textbf{面向海域自主侦测任务的边缘频谱智能处理方法}，构建频谱知识图谱、边缘在线学习与分布式可信决策一体化框架，实现端侧知识持续演化与多平台协同可信判决，为自主侦测提供敏捷可信的决策保障；最后，基于上述方法体系开展\textbf{海杂波场景电磁频谱数据智能处理原型系统验证}，推动关键技术落地实践，全面提升复杂海洋环境下电磁频谱的自主感知、轻量管理与智能决策能力（如图xxx所示）。

% 内容一
\textbf{研究内容一：海杂波干扰下微弱电磁频谱信号传播规律与微弱信号提取}

xxx

% A
\textbf{A. 复杂海洋环境下电磁频谱信号时空耦合规律建模}

复杂海洋环境下，海面粗糙散射、海杂波长尾起伏、低掠射角多径传播、海气边界层折射以及海警船观测平台运动共同作用，使嫌疑无人船、非法改装艇等弱目标电磁频谱信号在传播过程中呈现时变衰落、频谱展宽、空间相关和可观测性波动。现有方法多将海杂波作为背景噪声或独立干扰处理，难以解释海况、传播通道与目标信号显现程度之间的耦合关系，导致后续微弱信号提取缺乏稳定物理先验。本项目拟探索\textbf{复杂海洋环境下电磁频谱信号时空耦合传播规律建模}方法：首先，研究\textbf{海面散射--杂波统计联合表征方法}，刻画粗糙海面散射强度、海杂波幅度长尾、多普勒谱宽/频移、海尖峰率与风速、浪高、浪向、掠射角和极化方式之间的关联关系；其次，研究\textbf{低掠射角多径与海气折射耦合传播模型}，分析直达波、海面反射波和蒸发波导等异常传播路径对传播损耗、相位起伏和覆盖边界的影响；最后，研究\textbf{面向弱目标可观测性的时空传播规律提取方法}，构建传播损耗、杂波统计参数、短时谱结构和空间相关特征之间的联合描述，形成支撑微弱信号提取的场景先验、可感知性边界和关键传播参数。

% B
\textbf{B. 基于空域干扰对消的微弱目标频谱特征增强方法}

% C
\textbf{C. 基于多维联合表征的电磁频谱信号精确提取方法}
海上无人目标电磁辐射信号多数具有功率弱、持续时间短、频谱形态细碎、调制特征不完整和空间指向不稳定等特点，其有效信息分散在时域、频域、时频域、空域、能量域和调制域等多个观测维度中，单一谱域难以全面刻画目标信号的细粒度结构，导致复杂频谱中目标信号候选区域难以准确定位、目标信号与邻近频谱成分难以有效分离、关键特征参数难以稳定获取。本项目拟探索\textbf{基于多维联合表征的电磁频谱信号精确提取方法}：针对弱目标信号多域特征分散、不同观测维度难以统一描述的问题，研究\textbf{基于多域谱图张量编码的信号联合表征方法}，通过原始 IQ 片段、功率谱、时频谱、循环谱、调制谱和空间谱的同步构建，将时域波形结构、频域谱线分布、时频演化纹理、空域方位信息、能量变化规律和调制周期特征编码为统一的多域谱图张量，并结合时间配准、频率对齐、尺度归一和跨域特征关联建模，形成能够全面刻画目标信号多维属性的联合表征模型；针对复杂频谱中弱目标信号局部显著性不足、候选区域边界模糊、邻近频谱成分易与目标片段混叠的问题，提出\textbf{基于多证据一致性约束的候选区域检测与分离方法}，利用能量突变、谱线连续性、时频连通性、调制相关性和空域一致性等多类证据构建联合检测准则，通过多域显著性映射、候选区域连通跟踪和边界一致性约束，实现目标信号候选时频区域的快速定位，并进一步结合谱域结构约束和空调联合特征约束，分离目标信号片段与邻近频谱成分；针对目标信号精确提取过程中有效特征易损失、信号片段与特征参数难以同步恢复的问题，探讨\textbf{基于可微参数化重构的特征保持提取与参数反演方法}，将中心频率、带宽、持续时间、时频占用、能量包络、调制参数、循环特征和到达方位等关键参数嵌入可微信号重构模型，建立信号片段重构误差、特征保持约束和多域参数一致性之间的协同优化机制，实现无人目标电磁辐射信号片段、频谱边界、有效判别特征和关键参数的同步提取。



% 内容二
\textbf{研究内容二：有限算力端侧频谱数据的可信管理和分级存储方法}

xxx

% A
\textbf{A. 事件窗口驱动的频谱数据价值评估模型}

宽频段连续频谱监测会产生高吞吐、多冗余的时频数据，若全量输入 AI 模型精判并全量写入向量数据库，将显著增加海警船端推理开销和索引维护负担。本项目拟探索\textbf{事件窗口驱动的频谱数据价值评估模型}：针对连续频谱流中疑似无人艇电磁活动持续时间短、频点变化快、上下文证据易被固定切片割裂的问题，研究\textbf{基于流式变点检测的事件证据窗口构建方法}，通过在线背景建模、频谱状态突变感知和时频连通区域跟踪，将连续频谱流转化为包含触发上下文、频域邻域和时频演化轨迹的结构化事件证据单元；针对宽频段频谱事件难以全量进入 AI 大模型精判的问题，提出\textbf{基于 PU 风险估计的高召回事件识别方法}，利用少量已确认或高置信疑似无人艇事件与大规模未标注事件窗口学习候选事件评分函数，并通过候选排序机制优先识别疑似无人艇电磁事件和边界难判事件，从而压缩 AI 精判输入规模；针对事件窗口全量写入向量数据库易造成索引膨胀、相似检索冗余和查询效率下降的问题，探讨\textbf{基于子模覆盖增益的频谱事件向量入库判别方法}，结合已入库事件的表征分布、波形模板和频段覆盖情况，选择具有模式覆盖增益、稀缺波形补充和检索效用提升的事件窗口写入向量索引，从而在控制向量数据库规模的同时提升后续相似波形检索和目标复核效率。

% B
\textbf{B. 基于语义压缩的频谱向量数据库存储优化方法}

端侧环境存储空间高度受限，极大制约了海量电磁信号的连续感知与高效检索。本项目拟探索\textbf{基于语义压缩的频谱向量数据库存储优化方法}：针对原始频谱数据的物理驻留瓶颈，研究\textbf{面向强杂波干扰的语义提取与量化压缩方法}，通过提取核心物理语义并转化为紧凑离散码本，在严格控制误差的前提下实现海量数据的指数级压缩；针对图索引遍历引发的外存拥塞与I/O读放大问题，提出\textbf{面向异构介质的高局部性磁盘索引存储方法}，通过将跨页随机访问深度重构为局部顺序流，显著降低外存驻留图索引的访问开销；针对传统批量落盘更新占用额外内存且抢占底层检索带宽的问题，探讨\textbf{高频谱语义向量的无缓冲并发插入机制}，通过物理页内异地更新降低内存缓冲依赖，并有效平摊流式写入负载，实现动态高频写入场景下检索性能与索引维护的长效稳定。

% C
\textbf{C. 链上链下协同的可追溯频谱数据可信管理方法}

海域侦测频谱数据涉密且具监测取证与态势研判价值，要求数据不可篡改、全流程可追溯、来源可信，而端侧设备存储受限、无法部署完整区块链节点，原始IQ数据高吞吐、全量上链将带来难以承受的链上存储与带宽压力。本项目拟探索\textbf{链上链下协同的可追溯频谱数据可信管理方法}：针对频谱数据体量庞大、价值分布不均、全量上链开销过高的问题，研究\textbf{事件驱动的分层存证与动态摘要机制}，通过本体链下存储与摘要、元数据、操作日志链上存证的分层架构、事件窗口与数据价值驱动的多模动态摘要和适配海上弱网的异步上链协议，在最小化端侧算力与链上带宽开销的同时实现频谱数据的可信存证；针对多源频谱信息关联复杂、高价值证据查询定位与取证溯源困难的问题，提出\textbf{结合频谱特征的可验证分级索引构建方法}，利用按时段、频段、观测节点与信号样式的多维分级索引、查询结果与验证对象的同步生成机制和结构化检索与频谱相似性检索的双路融合，支持复杂条件下高价值频谱证据的快速定位与结果可验证；针对端侧弱算力、弱网下查询既要高效又要可信、且索引需随频谱数据高频变化持续维护的问题，探讨\textbf{基于学习索引的高效可验证查询与协同更新优化机制}，结合学习索引驱动的查询位置预测、基于链上摘要的正确性与完整性校验和索引与摘要的增量协同更新，在降低端侧维护开销的同时实现频谱数据全生命周期的可信追溯与防篡改。



\textbf{研究内容三：面向海域自主侦测任务的边缘频谱智能处理方法}

在海杂波干扰下的海域自主侦测任务中，前端微弱目标频谱提取结果和端侧可信数据管理结果只是形成自主侦测能力的基础，仍需进一步转化为边缘节点可理解、可更新、可协同使用的频谱认知与任务决策依据。该过程面临的关键挑战在于，海域边缘节点通常只能获得局部且不完备的频谱观测，单节点认知结果容易随环境扰动和任务状态变化而波动；多节点之间又存在观测视角、认知置信度和节点可信性差异，分散结果难以自然形成稳定一致的全局研判；同时，海上边缘节点受算力、带宽和时延约束，难以依赖全量数据集中汇聚或高交互复杂共识实现实时协同。鉴于此，本项目拟研究面向海域自主侦测任务的边缘频谱智能处理方法，在可信频谱数据和前端感知结果基础上，建立从频谱知识组织、单节点边缘认知到多节点可信协同决策的递进式处理机制，使频谱数据能够支撑边缘节点的即时理解、在线适应和分布式协同，形成面向复杂海域环境的可核验、可追溯、可执行自主侦测能力。


\textbf{A. 电磁频谱数据统一表征与知识库动态管理方法}

在海杂波干扰下的海域自主侦测任务中，前端感知结果和端侧可信数据若仅以分散记录形式存在，难以支撑边缘节点对当前频谱状态的稳定理解，也难以为后续在线学习和协同研判提供可复用的知识依据。其核心矛盾在于，频谱认知结果既受复杂海洋环境影响，又随任务过程持续演化，若缺乏统一表征和可信更新机制，历史经验、当前观测和协同反馈之间难以形成有效关联。本项目拟探索面向海杂波干扰环境的电磁频谱数据统一表征与知识库动态管理方法：针对频谱认知结果难以统一表达和跨场景复用的问题，研究\textbf{面向边缘认知的频谱数据统一表征方法}，将前端感知结果、端侧可信记录和协同反馈转化为可比较、可调用的统一知识表达，为边缘节点的认知判断和协同过程中的语义对齐提供基础；针对频谱知识难以从分散数据中沉淀为任务先验的问题，研究\textbf{频谱知识组织与任务关联建模方法}，围绕自主侦测任务建立频谱现象、环境影响和任务语义之间的关联关系，使知识库能够为单节点判断和多节点协同提供稳定先验；针对动态环境下知识更新容易受到偶发扰动和低可信反馈影响的问题，研究\textbf{可信反馈驱动的频谱知识库动态更新方法}，对进入知识库的新增信息进行可信筛选和状态管理，使稳定知识能够持续积累，待验证信息能够被审慎保留，从而形成面向海域自主侦测任务的可调用、可更新、可追溯频谱知识支撑能力。

\textbf{B. 频谱知识驱动的自主决策与边缘在线学习方法}

在海域边缘自主侦测任务中，单个边缘节点只能基于局部观测形成频谱认知结果，其判断容易受到海杂波干扰和观测条件变化的影响。若节点仅依据当前观测进行响应，局部误判可能被直接转化为任务行为；若缺乏面向动态环境的在线更新机制，节点又难以适应长期变化的海上频谱场景。本项目拟探索面向海杂波干扰环境的频谱知识驱动自主决策与边缘在线学习方法：针对单节点局部观测难以直接形成可靠任务语义的问题，研究\textbf{频谱知识约束的边缘频谱认知方法}，将目标频谱提取结果与频谱知识先验相结合，对当前观测结果进行可信解释，形成面向本地任务响应的边缘认知结果；针对边缘节点需要将认知结果转化为自主行为的问题，研究\textbf{面向自主侦测任务的单节点局部决策方法}，在知识先验和可信历史证据约束下建立局部决策机制，使节点能够依据当前认知状态生成与任务目标相匹配的本地响应策略；针对动态频谱环境下模型更新容易引入错误经验的问题，研究\textbf{可信样本驱动的边缘轻量在线学习方法}，对可用于更新的观测结果和反馈信息进行可信筛选，在有限算力条件下完成轻量自适应更新，使节点能够在不重复底层信号提取、不承担多节点确权决策的前提下，形成可信认知、局部自主决策和在线自适应能力。



\textbf{C. 频谱认知协同与分布式可信决策方法}

在海域多边缘节点自主侦测任务中，复杂海杂波、多径传播、同频干扰与链路受限等因素会导致各节点局部频谱认知结果存在不确定性、冲突性和可信性差异，单节点决策难以支撑稳定可靠的全局态势研判与协同任务执行。本项目拟探索面向海杂波干扰环境的频谱认知协同与分布式可信决策方法：针对多节点局部认知结果不一致和异常节点扰动问题，研究\textbf{基于低延时共识的决策结果一致性确权方法}，构建适配海上边缘低算力集群的轻量动态共识机制，通过节点动态可信评估与认知置信度驱动的加权共识，将存在差异的多节点局部认知和协同决策对象转化为可核验的一致性确权结果；针对多源异构观测下协同研判不稳定、融合决策缺乏可信约束的问题，研究\textbf{频谱知识驱动的分布式可信协同决策方法}，将频谱知识先验和确权可信约束嵌入多节点对齐、融合与推理过程，形成从频谱状态协同研判到任务级协同决策的可信决策链条，提升海杂波强干扰条件下协同决策的鲁棒性与可解释性；针对海上窄带链路下跨节点关键信息交互受限的问题，研究\textbf{有限链路的频谱感知结果轻量级安全传输方法}，以支撑共识确权和协同决策的轻量信息替代原始频谱流传输，并结合链下保存与链上锚定的可信回写机制，支撑协同感知、共识确权、融合决策和结果存证过程的可验证与可追溯，从而形成面向海域自主侦测任务的低时延、低开销、高鲁棒分布式可信协同决策能力。



% 内容四
\textbf{研究内容四：验证系统}

xxx
